\chapter{Data Collection and Processing}
\label{ch:data}

\section{Catalog and acquisition}

The catalog contains \CatalogPaperCount{} publication records derived from the TTLAB archive and stored in deterministic seed JSON. The discovery/import path normalises titles and URLs, assigns stable identifiers, preserves absent metadata as missing, and records provenance/review state. PDF acquisition uses allowed hosts and public IP resolution, bounded redirects, timeouts, content-length and byte limits, content-type and PDF-signature validation, and explicit failure categories. It does not bypass authentication, paywalls, or access controls.

The frozen audit records \AvailablePdfCount{} available PDFs, \NoPdfUrlCount{} records with no PDF URL, \PdfNetworkErrorCount{} network error, and \PdfNotFoundCount{} not-found result. Of the available documents, \EligiblePaperCount{} matched their catalog identity and entered the experimental corpus; \MismatchPaperCount{} possible PDF/title mismatches were excluded. The remaining \NoTextPaperCount{} catalog records are visible as needs-review/no-text records rather than being silently absent. \Cref{fig:ingestion} shows the full pipeline, and \cref{fig:corpus-status} shows the resulting paper-level eligibility categories.

\begin{figure}[htbp]
\centering
\resizebox{0.99\textwidth}{!}{\begin{tikzpicture}[node distance=8mm and 10mm]
  \node[flowbox] (archive) {TTLAB WordPress\\publication archive};
  \node[flowbox, right=of archive] (discover) {Parse metadata, URLs\\and stable paper ID};
  \node[datastore, right=of discover] (seed) {Seed JSON\\134 records};
  \node[flowbox, right=of seed] (import) {Upsert Paper and\\raw Author strings};

  \node[flowbox, below=14mm of import] (download) {Direct-PDF check,\\signature validation};
  \node[decision, left=12mm of download] (available) {PDF\\available?};
  \node[flowbox, left=12mm of available] (extract) {PyMuPDF page text +\\diagnostics + hash};
  \node[flowbox, left=of extract] (chunk) {Sentence/page chunks\\850 words, overlap 125};

  \node[datastore, below=14mm of chunk] (db) {SQLite\\Paper + Chunk};
  \node[flowbox, right=of db] (fts) {SQLite FTS5\\756 rows};
  \node[flowbox, right=of fts] (hash) {256-d feature hashing\\active index: 25 rows};
  \node[flowbox, right=of hash] (topics) {Rule-based topic and\\author aggregation};

  \draw[arrow] (archive) -- (discover);
  \draw[arrow] (discover) -- (seed);
  \draw[arrow] (seed) -- (import);
  \draw[arrow] (import) -- (download);
  \draw[arrow] (download) -- (available);
  \draw[arrow] (available) -- node[above,font=\scriptsize]{yes} (extract);
  \draw[arrow] (extract) -- (chunk);
  \draw[arrow] (chunk) -- (db);
  \draw[arrow] (db) -- (fts);
  \draw[arrow] (fts) -- (hash);
  \draw[arrow] (hash) -- (topics);
  \draw[dashedarrow] (available.south) -- ++(0,-8mm) -| node[pos=0.25,right,font=\scriptsize]{no: missing/failed status} (db.north);
\end{tikzpicture}
}
\caption[Deterministic ingestion and authoritative indexing pipeline.]{Deterministic ingestion and authoritative indexing pipeline. The optional OCR stage records status/provenance but was not used in the frozen corpus.}
\label{fig:ingestion}
\end{figure}

\begin{figure}[htbp]
\centering
\begin{tikzpicture}
\begin{axis}[
  ybar,
  width=0.92\linewidth,
  height=58mm,
  ymin=0,
  ymax=145,
  bar width=14pt,
  symbolic x coords={Eligible,No text,Mismatch},
  xtick=data,
  ylabel={Papers},
  nodes near coords,
  nodes near coords align={vertical},
  axis line style={draw=Slate},
  tick label style={font=\small},
  label style={font=\small},
  enlarge x limits=0.25,
  grid=major,
  grid style={gray!15}
]
\addplot[fill=Teal!65,draw=Teal!80!black] coordinates {(Eligible,96) (No text,36) (Mismatch,2)};
\end{axis}
\end{tikzpicture}

\caption[Frozen catalog eligibility.]{Frozen catalog eligibility: \EligiblePaperCount{} eligible full-text papers, \NoTextPaperCount{} records without eligible text, and \MismatchPaperCount{} PDF/title mismatch exclusions. These are inventory counts, not quality scores.}
\label{fig:corpus-status}
\end{figure}

\section{Extraction and OCR-aware diagnostics}

PyMuPDF extracts native text page by page and records page/character/word counts, empty pages, extraction status, and hashes. Content classification distinguishes textual, mixed, and unknown records. The repaired pipeline detects low-text/image-heavy pages and supports optional OCR with page-level engine/configuration, confidence/status, and review flags. OCR dependencies are optional and fail explicitly. All \CatalogPaperCount{} records in the frozen snapshot have \texttt{ocr\_status=not\_requested}; therefore, the thesis reports no OCR accuracy or latency.

The \AvailablePdfCount{} extracted PDFs contain \ExtractedPageCount{} pages and \ExtractedWordCount{} native-extraction words. These totals describe the local source snapshot; they are not included in the sanitized release. The suspected identity mismatches account for \MismatchChunkCount{} excluded chunks. The raw chunk table contains \RawChunkCount{} rows, while the eligible experiment and authoritative indexes contain \EligibleChunkCount{} rows.

\section{Page-aware chunking and sections}

Extracted sentences retain their page number. Chunks target 850 words with a 5,000-character limit and 125-word overlap; identifiers incorporate paper, sequence, page span, and source hash. Section names are inferred from headings and layout cues, with \texttt{Unknown} preserved when evidence is insufficient. Among eligible chunks, \EligibleUnknownChunkCount{} are unknown-section chunks (\EligibleUnknownChunkPercent\%).

A \SectionSilverCaseCount{}-case stratified AI-reviewed silver sample compared current section labels with source passages. Exact accuracy was \SectionSilverAccuracy{}, labeled coverage 0.900, macro precision \SectionSilverMacroPrecision{}, and macro recall \SectionSilverMacroRecall{}. Seven creation decisions required revision. These results support using section labels as a reranking feature with caution, not as authoritative document structure.

\section{Metadata and author identity}

The malformed \texttt{Click to View} author parser artifact was removed by an idempotent migration. The schema now stores canonical author rows, aliases, identity/merge state, and review status. The audit found \RawAuthorRowCount{} author rows: \CanonicalAuthorCount{} unresolved active canonical identities, two merged rows, and one invalid row. Exact alias collisions, eligible authorship mismatches, excluded-paper leakage, and prohibited expertise overclaims were all zero in the executed audit. \PossibleAuthorMergeCount{} possible initial/full-name or spelling pairs remain deliberately unmerged pending authoritative confirmation.

DOI, abstract, keyword, venue, and date fields carry provenance and review fields. Unverified metadata remain empty. Author-topic links are derived only from eligible indexed publication authorship and paper topics; their UI wording explicitly disclaims availability, endorsement, supervision, or expertise beyond this corpus.

\section{Corpus distribution and release boundary}

\begin{figure}[htbp]
\centering
\begin{tikzpicture}
\begin{axis}[
  ybar,
  width=0.97\linewidth,
  height=65mm,
  ymin=0,
  bar width=7pt,
  xlabel={Publication year in seed metadata},
  ylabel={Papers},
  xtick=data,
  x tick label style={rotate=45,anchor=east,font=\scriptsize},
  y tick label style={font=\small},
  label style={font=\small},
  grid=major,
  grid style={gray!15},
  nodes near coords,
  nodes near coords style={font=\tiny}
]
\addplot[fill=Navy!65,draw=Navy] table[x=year,y=count,col sep=comma] {generated/year_distribution.csv};
\end{axis}
\end{tikzpicture}

\caption[Publication-year distribution in catalog metadata.]{Publication-year distribution in the catalog metadata. Missing or uncertain years remain represented in the generated source data rather than being imputed.}
\label{fig:years}
\end{figure}

The year distribution in \cref{fig:years} characterises the seed, not full-text eligibility. Raw PDFs, extracted text, chunks, local database, and indexes are excluded from the sanitized release because TTLAB project authorization does not itself confer redistribution rights for every third-party paper. The release retains redistributable metadata, schemas, prompts without restricted passages, labels, aggregate/raw sanitized results, source hashes, and reacquisition instructions.
